\documentclass[11pt,a4paper]{article}
\usepackage[utf8]{inputenc}
\usepackage[T1]{fontenc}
\usepackage{lmodern}
\usepackage[margin=1in]{geometry}
\usepackage{amsmath,amssymb}
\usepackage{booktabs}
\usepackage{array}
\usepackage{xcolor}
\usepackage{tcolorbox}
\usepackage{enumitem}
\usepackage{fancyhdr}
\usepackage{hyperref}
\usepackage{graphicx}

\tcbuselibrary{skins,breakable}

% Colors
\definecolor{primary}{HTML}{1B365D}   % Navy
\definecolor{secondary}{HTML}{008080} % Teal
\definecolor{accent}{HTML}{C0392B}    % Crimson
\definecolor{boxbg}{HTML}{F8FAFC}     % Off-white
\definecolor{bordercolor}{HTML}{CBD5E1}

% Page styling
\pagestyle{fancy}
\fancyhf{}
\rhead{\textcolor{primary}{\small\textbf{Pipeline Document 02: Model Training}}}
\lhead{\small CUSAT MSc AI Mini-Project}
\cfoot{\thepage}
\renewcommand{\headrulewidth}{0.6pt}

\hypersetup{
    colorlinks=true,
    linkcolor=primary,
    urlcolor=secondary,
    citecolor=primary
}

\begin{document}

% Title Header
\begin{center}
    {\LARGE \textbf{\textcolor{primary}{02 $\cdot$ ECRFORMER \& PATCHGAN TRAINING PIPELINE}}}\\[0.4em]
    {\large \textbf{Generator Data Flow, Adversarial Critique, and High-Frequency Sharpening}}\\[0.6em]
    {\small \textbf{Project:} SAR-Guided Satellite Cloud Removal $\cdot$ ECRformer + Conditional PatchGAN}\\
    {\small \textbf{Authors:} Agnal K V $\cdot$ Bivin Balu $\cdot$ Vaishnav T Prakash $\quad|\quad$ \textbf{Supervisor:} Dr. Sabu M K}\\
    {\small Department of Computer Applications, Cochin University of Science and Technology (CUSAT)}
\end{center}
\vspace{-0.5em}
\noindent\rule{\linewidth}{1.2pt}

\vspace{0.8em}
\begin{tcolorbox}[colback=primary!5!white,colframe=primary,title=\textbf{Document Purpose \& Training Architecture},fonttitle=\bfseries]
This document covers the complete \textbf{model training pipeline}. It explains how 7-channel multi-modal inputs ($S_1$ SAR, $S_{2,\text{cloudy}}$ optical, and soft mask $M$) pass through the \textbf{ECRformer Generator} ($G_\theta$), how the \textbf{11-channel Conditional PatchGAN Critic} ($D_\phi$) evaluates local texture fidelity, how \textbf{Opaque Cloud Cutout} prevents modality cheating, and \textbf{why these methods were chosen over standard alternatives}.
\end{tcolorbox}

\section{The Two-Player Adversarial Architecture}

Training operates as a minimax game between two specialized networks:
\begin{enumerate}[leftmargin=1.5em, itemsep=0.3em]
    \item \textbf{Generator $G_\theta$ (ECRformer-Lite, 3.62M params):} Ingests 7-channel multi-modal stack and predicts unblended clear reflectance $P_{\text{dir}} \in \mathbb{R}^{B \times 4 \times 128 \times 128}$.
    \item \textbf{Discriminator $D_\phi$ (11-channel PatchGAN, 2.77M params):} Ingests 11-channel conditional stack and evaluates local $16 \times 16$ spatial patch authenticity.
\end{enumerate}

\section{Forward Data Flow Through ECRformer}

\begin{figure}[h!]
\centering
\small
\renewcommand{\arraystretch}{1.25}
\resizebox{\linewidth}{!}{%
\begin{tabular}{lllll}
\toprule
\textbf{Stage} & \textbf{Operation / Layer} & \textbf{Input Shape} & \textbf{Output Shape} & \textbf{Core Role in Cloud Removal} \\
\midrule
\textbf{0. Input Split} & Slicing across channel axis & $[B, 7, 128, 128]$ & $S_2:[B,4], S_1:[B,2], M:[B,1]$ & Isolates optical, SAR radar, and cloud mask \\
\textbf{1. Decoupled Stem} & Two separate $7\times 7$ Convs & $S_2:[B,4], S_1:[B,2]$ & $F_{\text{opt}}:[B,32], F_{\text{sar}}:[B,32]$ & Domain-specific low-level feature extraction \\
\textbf{2. 4-Stage Encoder} & Depth $[2, 2, 1, 1]$ & $[B, 32, 128, 128]$ & $[B, 256, 16, 16]$ & Multi-scale Cross-Modality Attention (CMA) \\
\textbf{3. TSA Bottleneck}& Time-Space Attention & $[B, 256, 16, 16]$ & $[B, 256, 16, 16]$ & Global non-local spatial context modeling \\
\textbf{4. Symmetric Decoder}& Transposed Convs + CBAM & $[B, 256, 16, 16]$ & $[B, 32, 128, 128]$ & Upsampling with filtered encoder skip connections \\
\textbf{5. Residual Head} & $3\times 3$ Conv + Mean Prior & $[B, 32, 128, 128]$ & $[B, 4, 128, 128]$ & Outputs unblended direct prediction $P_{\text{dir}}$ \\
\bottomrule
\end{tabular}%
}
\caption{Layer-by-layer tensor progression throughout ECRformer Generator.}
\end{figure}

\subsection{Cross-Modality Attention (CMA)}
Inside each encoder stage, optical queries ($Q$) query radar keys and values ($K, V$):
\begin{equation}
\text{CMA}(Q, K, V) = \text{Softmax}\left( \frac{Q K^T}{\sqrt{d_k}} \right) V
\end{equation}
This allows cloudy optical regions to dynamically import physical dielectric surface contours from the penetrating SAR signal.

\newpage

\section{PatchGAN Critique \& Opaque Cloud Cutout}

\subsection{Overcoming the "Blurriness Trap"}
Standard $L_1$ and $L_2$ losses predict the conditional mean $\mathbb{E}[Y | X]$, which turns highways and building blocks into smooth grey smudges. PatchGAN solves this by critiquing high frequencies across $16 \times 16$ local patches.

\subsection{Opaque Cloud Cutout ($65\%$ Modality Dropout)}
In thin haze, optical ground lines faintly leak through $S_{2,\text{cloudy}}$. If left unchecked, the model cheats by ignoring SAR. In $65\%$ of training crops, we erase optical ground edges under the mask using 2-pass $15 \times 15$ pooling and bright cloud veiling, \textbf{forcing the network to learn SAR feature reconstruction}.

\subsection{Composite Training Loss}
\begin{equation}
\mathcal{L}_{\text{total}} = \mathcal{L}_{L1} + \lambda_{\text{SDFL}} \mathcal{L}_{\text{SDFL}} + w_{\text{adv}} \mathcal{L}_{\text{adv}} + w_{\text{FM}} \mathcal{L}_{\text{FM}} + \lambda_{\text{corr}} \mathcal{L}_{\text{corr}}
\end{equation}
with $w_{\text{adv}} = 0.04$ and $w_{\text{FM}} = 0.35$.

\section{Why We Used These Methods Over Others}

\begin{tcolorbox}[colback=white,colframe=primary,title=\textbf{1. Why ECRformer over Standard Convolutional U-Nets (e.g. Pix2Pix, DSen2-CR)?}]
\textbf{Alternative:} Standard U-Net with local $3 \times 3$ convolutions.
\begin{itemize}[leftmargin=1.2em, itemsep=0.2em]
    \item \textbf{Failure of Alternative:} Local convolutions cannot see beyond $5 \times 5$ pixels. Under vast regional cloud cover ($>50\text{ km}$), convolutions see only white pixels and cannot borrow context from distant clear ground.
    \item \textbf{Why Our Choice Wins:} Self-attention gives ECRformer an infinite receptive field, enabling it to reference clear ground context $10+\text{ km}$ away.
\end{itemize}
\end{tcolorbox}

\begin{tcolorbox}[colback=white,colframe=secondary,title=\textbf{2. Why an 11-Channel Multi-Modal Discriminator over Optical-Only Critics?}]
\textbf{Alternative:} Feeding only the generated optical image $P_{\text{dir}}$ to the discriminator.
\begin{itemize}[leftmargin=1.2em, itemsep=0.2em]
    \item \textbf{Failure of Alternative:} Optical-only critics only verify whether an image looks like Earth terrain; the generator can hallucinate plausible farms where the ground is actually an industrial port.
    \item \textbf{Why Our Choice Wins:} Conditioning on Sentinel-1 SAR radar ($11\text{ channels}$) forces the generated optical textures to strictly obey the physical microwave backscatter of the real ground below.
\end{itemize}
\end{tcolorbox}

\begin{tcolorbox}[colback=white,colframe=primary,title=\textbf{3. Why Scale-Decoupled Frequency Loss (SDFL) over Standard $L_1$ Alone?}]
\textbf{Alternative:} Mean Absolute Error ($L_1$) across pixels.
\begin{itemize}[leftmargin=1.2em, itemsep=0.2em]
    \item \textbf{Failure of Alternative:} $L_1$ is dominated by low-frequency illumination, rewarding blurry mean predictions.
    \item \textbf{Why Our Choice Wins:} SDFL uses a Laplacian pyramid to isolate high-frequency octaves, penalizing edge gradient errors independently and boosting Detail Ratio from $0.354$ to $0.931$ (+142\% sharper).
\end{itemize}
\end{tcolorbox}

% Quick Viva Sheet
\section{15-Second Viva Examiner Defense Quick-Sheet}

\begin{description}[leftmargin=1.5em, style=nextline]
    \item[\textbf{Q: How many parameters are trained in your system?}]
    \textit{Answer:} The generator contains $3.62$ Million parameters ($\approx 14\text{ MB}$), while the PatchGAN critic contains $2.77$ Million parameters, enabling efficient training on a single GPU.
    
    \item[\textbf{Q: What is the purpose of the Opaque Cloud Cutout augmentation?}]
    \textit{Answer:} It erases leaked optical ground textures in $65\%$ of training crops to prevent the network from taking optical shortcuts, forcing it to utilize Sentinel-1 SAR radar.
\end{description}

\end{document}
